% Fiche d'examen -- Optimisation (CentraleSupelec, cours G. Allaire / V. Lescarret) -- VERSION IMPRIMEE
% A4 recto-verso. Contenu : fiche_contenu.tex, tire uniquement des slides de Opti_CS/Slides_Cours_officiel.
% Compilation : pdflatex fiche_opti.tex (une passe suffit).
\documentclass[a4paper]{article}
\usepackage[utf8]{inputenc}
\usepackage[T1]{fontenc}
\usepackage{lmodern}
\usepackage[a4paper,margin=4mm,top=3.5mm,bottom=3.5mm]{geometry}
\usepackage{amsmath,amssymb,mathtools}
\usepackage{multicol}
\usepackage[dvipsnames]{xcolor}
\usepackage{tikz}
\usetikzlibrary{shapes.geometric}
\usepackage{ragged2e}
\usepackage{microtype}
\pagestyle{empty}
\setlength{\parindent}{0pt}
\setlength{\parskip}{0pt}
\setlength{\columnsep}{2.4mm}
\setlength{\columnseprule}{0.2pt}
\setlength{\multicolsep}{0pt}
\setlength{\abovedisplayskip}{0pt}\setlength{\belowdisplayskip}{0pt}
\setlength{\abovedisplayshortskip}{0pt}\setlength{\belowdisplayshortskip}{0pt}
\setlength{\arraycolsep}{1.5pt}
\thickmuskip=2.5mu plus 2mu \medmuskip=2mu plus 1mu \thinmuskip=1.5mu

% ---------- notations ----------
\newcommand{\R}{\mathbb{R}}\newcommand{\N}{\mathbb{N}}\newcommand{\Z}{\mathbb{Z}}
\newcommand{\sca}[2]{\langle #1,#2\rangle}
\newcommand{\nr}[1]{\|#1\|}
\newcommand{\eps}{\varepsilon}
\newcommand{\Lg}{\mathcal{L}}

% ---------- couleurs ----------
\definecolor{cT}{HTML}{1F4E79}  % titres et noms de resultats
\definecolor{cP}{HTML}{C0392B}  % avertissements
\definecolor{cK}{HTML}{1E7B34}  % CN / CS / methode

% ---------- structure d'un enonce ----------
\newcommand{\B}{\par\noindent\hangindent=0.8em\hangafter=1\makebox[0.8em][l]{\textbullet}}
\newcommand{\Th}[1]{\textbf{\color{cT}#1}}
\newcommand{\Hyp}{\textbf{Hyp.}~}
\newcommand{\Alors}{$\boldsymbol{\Rightarrow}$~}
\newcommand{\CN}{\textbf{\color{cK}[CN]}~}
\newcommand{\CS}{\textbf{\color{cK}[CS]}~}
\newcommand{\Mt}{\textbf{\color{cK}$\blacktriangleright$}\,}
\newcommand{\Ex}{\textit{Ex.}~}
\newcommand{\pg}[1]{{\color{cP}#1}}
% precision absente des slides (correction d'un enonce de slide incomplet ou faux)
\newcommand{\warnsym}{\tikz[baseline=-0.6ex]\node[draw=cP,line width=0.45pt,regular polygon,regular polygon sides=3,inner sep=0pt,minimum size=1.05em,text=cP,font=\fontsize{5.6}{6}\selectfont\bfseries]{!};}
\newcommand{\Warn}[1]{{\color{cP}\warnsym\if\relax\detokenize{#1}\relax\else\,(#1)\fi}}
% titres de section : bandeau de couleur
\newcommand{\Sec}[1]{\par\vspace{1.4pt}\noindent\colorbox{cT}{\parbox{\dimexpr\linewidth-2\fboxsep}{\color{white}\bfseries #1}}\par\vspace{0.8pt}}
\newcommand{\TitreFiche}{\fontsize{10.5}{11.5}\selectfont\bfseries\color{cT}}
\newcommand{\Legende}{\fontsize{7}{8}\selectfont}
\begin{document}
\fontsize{9.1}{10.24}\selectfont
% Contenu de la fiche d'examen, commun aux deux habillages :
%   fiche_opti.tex (imprimee, pdflatex) et fiche_opti_manuscrite.tex (manuscrite, lualatex).
% Source : uniquement les slides du dossier Opti_CS/Slides_Cours_officiel
%   (SlidesAllaire, ProgLin, t24_cours5 [EDP], note Fonctions-implicites).
% \Warn{...} signale une precision qui n'est PAS dans les slides : elle corrige un enonce
%   incomplet ou faux des slides, sans rien y ajouter d'autre.

\noindent{\TitreFiche OPTIMISATION -- fiche d'examen}\hfill{\Legende CN = condition nécessaire, CS = suffisante \ \ \Warn{} = précision absente des slides}\par
\vspace{2pt}
\begin{multicols}{3}\RaggedRight

%=====================================================================
\Sec{1. Cadre et vocabulaire}
\B Trouver $u\in X_{ad}\subset X$ tel que $J(u)\le J(v)$ $\forall v\in X_{ad}$ ; $X$ Hilbert ($\R^n$, espace fonctionnel) ou ensemble dénombrable ; $X_{ad}=K$ \textbf{fermé} (états admissibles : contraintes d'égalité, d'inégalité) ; $J$ critère (coût).
\B $\inf_KJ$ si on ne sait pas si $J$ admet un min, $\min_KJ$ sinon.
\B $u$ \textbf{min local} de $J$ sur $K$ : $\exists\delta>0$, $\forall v\in K$, $\nr{u-v}\le\delta\Rightarrow J(u)\le J(v)$.
\B \textbf{Suite minimisante} : $u_n\in K$, $J(u_n)\to\inf_KJ$ ; existe toujours (déf. de l'inf).
\B \Ex Markowitz : $\min\frac12\,{}^t\!XVX$ s.c. $\sum x_i=1$, $x_i\ge0$ (ou $\min-{}^t\!RX$ s.c. $W(X)\le\rho$) ; voyageur de commerce : permutation $\sigma$ minimisant $\sum d(M_{\sigma(i)},M_{\sigma(i+1)})$ (discret) ; identification : $\min\nr{w(\theta)-w_{obs}}$ s.c. $M(u,\theta,v)=0$.

%=====================================================================
\Sec{2. Existence et unicité}
\B \textbf{Déf.} (dim finie) $J$ \textbf{coercive} : $\lim_{|x|\to\infty}J(x)=\infty$. $J$ \textbf{sci} : $\forall(u_n)$ suite de $K$, $u_n\to u\Rightarrow\liminf J(u_n)\ge J(u)$.
\B \Th{Th 2.2.1} \Hyp $K\subset\R^n$ \textbf{fermé}, $J$ sci et coercive sur $K$. \Alors $J$ a au moins un point de min, et de toute suite minimisante on extrait une sous-suite qui converge vers un tel point.
\B \Th{Prop 2.3.4} \Hyp $K$ convexe, $J$ convexe sur $K$. \Alors tout min local est \textbf{global} ; l'ensemble des points de min est convexe (éventuellement vide). \textbf{Si de plus} $J$ strictement convexe : \textbf{au plus un} point de min.
\B Convexes : $J_1,J_2$ convexes $\Rightarrow J_1+J_2$ et $\max(J_1,J_2)$ convexes ; $J$ convexe est continue \Warn{en dim. finie, sur un ouvert} ; son graphe est convexe et au-dessus de ses tangentes.
\B \Th{Th 2.3.9} \Hyp $K$ convexe fermé non vide d'un Hilbert $V$, $J$ convexe, sci et coercive sur $K$. \Alors il existe un min de $J$ sur $K$ (unicité non affirmée).
\B \textbf{Déf.} $J$ \textbf{$\alpha$-convexe} (fortement convexe) sur $K$ convexe : $\exists\alpha>0$, $\forall u,v\in K$, $\forall\theta\in[0,1]$, $J(\theta u+(1-\theta)v)\le\theta J(u)+(1-\theta)J(v)-\alpha\frac{\theta(1-\theta)}2\nr{u-v}^2$.
\B \Th{Th 2.3.8} \Hyp $K$ convexe fermé non vide d'un Hilbert $V$, $J$ \textbf{$\alpha$-convexe continue} sur $K$. \Alors il existe un \textbf{unique} min $u$, et $\forall v\in K$, $\nr{v-u}^2\le\frac4\alpha\big(J(v)-J(u)\big)$ ; toute suite minimisante converge vers $u$.

%=====================================================================
\Sec{3. Dérivées}
\B \Th{Déf 2.4.1} (Fréchet) $J$ dérivable en $u$ s'il existe $L\in V'$ (forme linéaire continue) : $J(u+w)=J(u)+L(w)+o(w)$, $\frac{|o(w)|}{\nr w}\to0$. $L=J'(u)$. Riesz : $L(v)=\sca pv$, on note $p=J'(u)$. Dans $\R^n$ : $\partial_iJ(u)=L(e_i)$.
\B \Th{Déf 2.4.7} $J$ deux fois dérivable en $u$ : $J$ dérivable au voisinage de $u$ et $J'$ dérivable en $u$ : $J'(u+w)=J'(u)+J''(u,w)+o(w)$ dans $V'$ ; $J''(u)(w,v):=\sca{J''(u,w)}v$.
\B \Th{Lemme 2.4.8} $J$ deux fois dérivable $\Rightarrow J(u+w)=J(u)+\sca{J'(u)}w+\frac12J''(u)(w,w)+o(\nr w^2)$.

%=====================================================================
\Sec{4. Conditions d'optimalité}
\B \Th{Th 2.5.1} (inéquation d'Euler) \Hyp $K$ \textbf{convexe}, $u\in K$, $J$ dérivable en $u$.\\
\CN \textbf{Si} $u$ \textbf{min local} de $J$ sur $K$ \Alors $\sca{J'(u)}{v-u}\ge0$ $\forall v\in K$.\\
\CS \textbf{Si} $u\in K$ vérifie cette inégalité \textbf{et} $J$ convexe \Alors $u$ min \textbf{global} sur $K$.
\B \Th{Déf 2.5.4} Cône des directions admissibles en $u\in K$ : $K(u)=\{w\in V:\exists u_n\in K\to u,\ \exists\eps_n\searrow0,\ \frac{u_n-u}{\eps_n}\to w\}$.
\B \Th{Prop 2.5.5} \CN \Hyp $J$ dérivable en $u$. \textbf{Si} $u$ \textbf{min local} de $J$ sur $K$ \Alors $\sca{J'(u)}w\ge0$ $\forall w\in K(u)$.

%=====================================================================
\Sec{5. Égalités : Euler-Lagrange}
\B $K=\{v\in V : F_1(v)=\dots=F_M(v)=0\}$. Si les $F_i$ sont dérivables en $u$ : $K(u)\subset\bigcap_i(F_i'(u))^\perp$.
\B \Th{Th 2.5.6} \CN \Hyp $u\in K$, $J$ dérivable en $u$, $F_i$ $C^1$ au voisinage de $u$, $\big(F_i'(u)\big)$ \textbf{linéairement indépendants} (cas régulier). \textbf{Si} $u$ \textbf{min local} de $J$ sur $K$ \Alors $\exists\lambda_1,\dots,\lambda_M\in\R$ (multiplicateurs de Lagrange) : $J'(u)+\sum\lambda_iF_i'(u)=0$.
\B \Th{Prop 2.5.13} \CN d'ordre 2. \Hyp celles du Th 2.5.6 $+$ $J$, $F_i$ deux fois continûment dérivables au voisinage de $u$ ; $\lambda$ = multiplicateurs du Th 2.5.6. \textbf{Si} $u$ \textbf{min local} de $J$ sur $K$ \Alors $\big(J''(u)+\sum\lambda_iF_i''(u)\big)(w,w)\ge0$ $\forall w\in K(u)$.
\B \textbf{Fonctions implicites} (preuve, chemins admissibles) : $F(u(t))=0$, $u(0)=u$ ; $F'(u)\ne0\Rightarrow\exists n$, $\partial_nF(u)\ne0$ ; si $n=1$, $u=(u_1,\tilde u)$ et $u_1=f(\tilde u)$ ; tout chemin $\tilde u=\tilde U(t)$ donne $u(t)=\big(f(\tilde U(t)),\tilde U(t)\big)$, chemin de $F^{-1}(0)$.

%=====================================================================
\Sec{6. Inégalités : KKT}
\B $K=\{v\in V : F_1(v)\le0,\dots,F_M(v)\le0\}$. \Th{Déf 2.5.14} Contraintes \textbf{actives} : $I(u)=\{i:F_i(u)=0\}$.
\B Si $I(u)\ne\emptyset$ : $K(u)$ = cône tangent des directions rentrantes, $K(u)\subset\{d:\sca{F_i'(u)}d\le0,\ i\in I(u)\}$.
\B \Th{Déf 2.5.15} Contraintes \textbf{qualifiées} en $u\in K$ $\Leftrightarrow$ $\exists\bar w\in V$, $\forall i\in I(u)$ : $\sca{F_i'(u)}{\bar w}<0$ si $F_i$ non linéaire, $\sca{F_i'(u)}{\bar w}=0$ si $F_i$ affine.
\B \Th{Lemme} \Hyp $I(u)\ne\emptyset$, $\big(F_i'(u)\big)_{i\in I(u)}$ \textbf{libre}. \Alors contraintes qualifiées en $u$.
\B \Th{Th 2.5.17} (KKT) \CN \Hyp $J$, $F_i$ dérivables en $u$, contraintes \textbf{qualifiées} en $u$. \textbf{Si} $u$ \textbf{min local} de $J$ sur $K$ \Alors $\exists\lambda_1,\dots,\lambda_M\ge0$ : $J'(u)+\sum\lambda_iF_i'(u)=0$ et $\lambda_i=0$ si $F_i(u)<0$.
\B \Th{Th 2.5.24} (KKT généralisé) \CN $K=\{G_1=\dots=G_N=0,\ F_1\le0,\dots,F_M\le0\}$. \Hyp $J,G_i,F_j$ dérivables en $u$, $\big(G_i'(u)\big)_{i\le N}$ libre, $\exists\bar w\in\bigcap_i(G_i'(u))^\perp$ : $\sca{F_j'(u)}{\bar w}<0$. \textbf{Si} $u$ \textbf{min local} \Alors $\exists\mu\in\R^N$, $\lambda\in(\R^+)^M$ : $J'(u)+\sum\mu_iG_i'(u)+\sum\lambda_iF_i'(u)=0$, $\lambda_i=0$ si $F_i(u)<0$.

%=====================================================================
\Sec{7. Lagrangien et dualité}
\B \Th{Déf 2.5.50} $\Lg(v,\mu)=J(v)+\sum\mu_iF_i(v)$, $(v,\mu)\in V\times(\R^+)^M$.
\B \Th{Lemme 2.5.21} $\displaystyle\inf_{F_i(v)\le0}J(v)=\inf_{v\in V}\sup_{\mu\in(\R^+)^M}\Lg(v,\mu)$.
\B \Th{Déf 2.6.7} $I(v)=\sup_{\mu\ge0}\Lg(v,\mu)$, $G(\mu)=\inf_{v\in V}\Lg(v,\mu)$. \textbf{Primal} : $\inf_VI$ (c'est le problème initial $\inf_{F\le0}J$). \textbf{Dual} : $\sup_{\mu\ge0}G$. Toujours $\inf I\ge\sup G$ ; la différence est le \textbf{saut de dualité}.
\B \Th{Th 2.6.9} (dualité convexe) \Hyp $J$, $F_i$ \textbf{convexes}, $u$ solution du primal. \Alors $\exists\lambda\in(\R^+)^M$ : $(u,\lambda)$ \textbf{point-selle} de $\Lg$ sur $V\times(\R^+)^M$ et $I(u)=\min_vI(v)=\max_\mu G(\mu)=G(\lambda)$. \Warn{il faut aussi les contraintes qualifiées en $u$ ; sinon faux : $\min x$ s.c. $x^2\le0$, $\sup G=0$ non atteint} (point-selle : $\Lg(u,\mu)\le\Lg(u,\lambda)\le\Lg(v,\lambda)$ $\forall v,\mu$).

%=====================================================================
\Sec{8. Algorithmes sans contraintes}
\B \textbf{Gradient à pas fixe} : $u_{n+1}=u_n-\mu J'(u_n)$, $u_0\in V$.
\B \Th{Th 3.2.6} \Hyp $J$ \textbf{$\alpha$-convexe} différentiable, $J'$ \textbf{$L$-lipschitzien} sur $V$ ($\nr{J'(u)-J'(v)}\le L\nr{u-v}$), $0<\mu<2\alpha/L^2$. \Alors $\forall u_0$, $u_n\to u$ (le min) et $\nr{u_n-u}\le\gamma^n\nr{u_0-u}$, $\gamma=\sqrt{1-2\alpha\mu+\mu^2L^2}<1$.
\B \Th{Déf 3.3.19} $J$ convexe sur $\R^n$ : \textbf{sous-différentiel} $\partial J(x)=\{p:J(y)-J(x)\ge p\cdot(y-x)\ \forall y\}$ (sous-gradients). $J$ dérivable en $x\Rightarrow\partial J(x)=\{J'(x)\}$ ; $0\in\partial J(x)\Rightarrow x$ min de $J$ ; $J=\sup_\lambda J_\lambda$ ($J_\lambda$ convexes dérivables) : si $J(x)=J_\lambda(x)$, alors $J_\lambda'(x)\in\partial J(x)$.
\B \textbf{Sous-gradient} : $x_{n+1}=x_n-\rho_n\frac{p_n}{\nr{p_n}}$, $p_n\in\partial J(x_n)$. \Th{Prop 3.3.23} \Hyp $J$ convexe, coercive, $\partial J(x)\ne\emptyset$ partout, localement lipschitzienne ($\forall M$, $\exists L_M$ : $\nr x+\nr y\le M\Rightarrow|J(x)-J(y)|\le L_M\nr{x-y}$), $(\rho_n)\in\ell^2\setminus\ell^1$. \Alors l'algorithme converge.
\B \textbf{Newton} : racine de $F:\R^n\to\R^n$ : $u_{n+1}-u_n=-(F'(u_n))^{-1}F(u_n)$, $u_0$ proche de $u^*$, $F'(u^*)$ inversible ; minimisation ($J$ $C^2$) : $u_{n+1}-u_n=-(J''(u_n))^{-1}J'(u_n)$, $u_0$ proche d'un point critique $u^*$, $J''(u^*)$ inversible.
\B \Th{Prop 3.5.1} \Hyp $F$ $C^2$ de $\R^n$ dans $\R^n$, $u^*$ zéro \textbf{régulier} ($F(u^*)=0$, $F'(u^*)$ inversible). \Alors $\exists\eps>0$, $0<C<1/\eps$ : si $\nr{u^*-u_0}\le\eps$, convergence \textbf{quadratique} $\nr{u_{n+1}-u^*}\le C\nr{u_n-u^*}^2$ ; $C\le\sup_{|v|=1}\nr{(F'(u^*))^{-1}F''(u^*)(v)}$.

%=====================================================================
\Sec{9. Algorithmes sous contraintes}
\B \textbf{Newton, égalités} $\min_{G(v)=0}J$ : un min est point critique du Lagrangien, $\nabla\Lg(u,\lambda)=\big(J'(u)+\lambda G'(u),\,G(u)\big)\in V\times\R^M$ ; $\Lg''(u,\lambda)(v,p;w,q)=J''(u)(v,w)+\lambda\cdot G''(u)(v,w)+q\cdot\sca{G'(u)}v+p\cdot\sca{G'(u)}w$ ; si $V=\R^n$ : $\Lg''=\begin{psmallmatrix}J''(u)+\lambda G''(u)&{}^tG'(u)\\G'(u)&0\end{psmallmatrix}$, et $\Lg''\ge0\Leftrightarrow J''+\lambda\cdot G''\ge0$ sur $\bigcap_j(G_j'(u))^\perp$ \Warn{positivité testée sur les $(v,p)$ avec $G'(u)v=0$}.
\B \textbf{Gradient projeté} ($K$ convexe fermé) : $u_{n+1}=P_K\big(u_n-\mu J'(u_n)\big)$. \Th{Th 3.4.1} \Hyp $J$ $\alpha$-convexe différentiable, $J'$ $L$-lipschitzien, $0<\mu<2\alpha/L^2$. \Alors l'algorithme converge.
\B \textbf{Uzawa} (inégalités) : $\Lg(v,\mu)=J(v)+\mu\cdot F(v)$, $\eps>0$, $\mu_0\in\R^M_+$ ; $u_n$ : $\Lg(u_n,\mu_n)=\inf_V\Lg(v,\mu_n)$, puis $\mu_{n+1}=P_{\R^M_+}\big(\mu_n+\eps F(u_n)\big)$. \Th{Th 3.4.3} \Hyp $J$ $\alpha$-convexe différentiable, $F$ convexe et $L$-lipschitzienne de $V$ dans $\R^M$, $0<\eps<2\alpha/L^2$. \Alors Uzawa converge.
\B \textbf{Pénalisation} : $J_\eps(v)=J(v)+\frac1\eps\sum\max(0,F_i(v))^2$ sur $\R^n$. \Th{Prop 3.4.6} \Hyp $J$ continue, strictement convexe, coercive ; $F_i$ convexes continues ; $K=\{F\le0\}\ne\emptyset$ ; $u$ unique solution, $u_\eps$ min de $J_\eps$. \Alors $u_\eps\to u$ quand $\eps\to0$.

%=====================================================================
\Sec{10. EDP : espaces de Sobolev}
\B $\partial^\alpha_xu=\partial^{\alpha_1}_{x_1}\cdots\partial^{\alpha_d}_{x_d}u$ (au sens des distributions). $\Omega$ ouvert de $\R^d$ : $H^n(\Omega)=\{u\in L^2:\partial^\alpha u\in L^2,\ |\alpha|_1\le n\}$ ; $H^1(\Omega)=\{v\in L^2:\nabla v\in(L^2)^d\}$.
\B \Th{Th} $H^n(\Omega)$ muni de $(u,v)_{H^n}=\sum_{|\alpha|\le n}\int_\Omega\partial^\alpha u\,\partial^\alpha v$ est un \textbf{Hilbert} (preuve : $v_n$ de Cauchy $\Rightarrow v_n\to v$, $v_n'\to w$ dans $L^2$, et $\langle T_v',\varphi\rangle=-\lim\int v_n\varphi'=\lim\int v_n'\varphi=\int w\varphi$, donc $v'=w$).
\B \Th{Th} (dim 1) $I$ intervalle ouvert \textbf{borné} : $H^1(I)\subset C^0(I)$ et $\exists C(I)$, $\nr u_\infty\le C\nr u_{H^1(I)}$ (th. fondamental de l'analyse sur $[y,x]$, intégré en $y$).
\B Dim $\ge2$ : en général $H^1(\Omega)\not\subset C(\Omega)$ : $|\ln(x^2+y^2)|^{1/3}\in H^1(D(0,1))$, non continue. $\Omega$ borné \textbf{régulier} : chaque point du bord a un (hyper)plan tangent, i.e. une seule normale. $k>d/2$, $\Omega$ régulier $\Rightarrow H^k(\Omega)\subset C^0(\bar\Omega)$.
\B \textbf{Trace} $\gamma_0:u\mapsto u_{|\partial\Omega}$ (linéaire). $I=]a,b[$ : $\max(|u(a)|,|u(b)|)\le c\nr u_{H^1(I)}$. \Th{Th} $\Omega$ borné régulier $\Rightarrow\gamma_0u$ définie p.p., $\in L^2(\partial\Omega)$, $\nr{\gamma_0u}_{L^2(\partial\Omega)}\le C_\Omega\nr u_{H^1(\Omega)}$.
\B \Th{Green} ($\Omega$ borné régulier) $w,v\in H^1$ : $\int_\Omega w\,\partial_iv=-\int_\Omega v\,\partial_iw+\int_{\partial\Omega}wv\,n_i$. $u\in H^2$, $v\in H^1$ : $\int_\Omega\nabla u\cdot\nabla v+\int_\Omega(\Delta u)v=\int_{\partial\Omega}v\,\nabla u\cdot n$.
\B \textbf{$H^1_0$} ($\Omega$ borné régulier) : $\mathcal D(\Omega)$ dense dans $L^2$, pas dans $H^1$ ; $H^1_0(\Omega)=\overline{\mathcal D(\Omega)}^{\,H^1}$. $H^1_0\subset H^1$, $H^1_0\ne H^1$ ; $H^1_0=\gamma_0^{-1}(\{0\})$ ; $(H^1_0,\nr\cdot_{H^1})$ Hilbert.
\B \Th{Th (Poincaré)} \Hyp $\Omega$ ouvert \textbf{borné} (ou borné dans une direction). \Alors $\exists k_\Omega$, $\forall v\in H^1_0(\Omega)$, $\nr v_{L^2}\le k_\Omega\nr{\nabla v}_{L^2}$. Faux pour $v\in H^1\setminus H^1_0$ (non nul au bord).
\B \Th{Th} $(u,v)\mapsto\int_\Omega\nabla u\cdot\nabla v$ est un produit scalaire sur $H^1_0$ (Hilbert) ; $|v|_{H^1}=\nr{\nabla v}_{L^2}$ est une norme sur $H^1_0$, équivalente à $\nr\cdot_{H^1}$ : l'outil des problèmes de Dirichlet homogènes.

%=====================================================================
\Sec{11. EDP : Lax-Milgram, 5 étapes}
\B $H'=\mathcal L(H,\R)$. \Th{Riesz-Fréchet} $\forall L\in H'$, $\exists!f\in H$ : $L(v)=\sca fv_H$ $\forall v$.
\B $a$ bilinéaire \textbf{continue} : $|a(u,v)|\le C\nr u\nr v$ ; \textbf{coercive} : $\exists\alpha>0$, $a(u,u)\ge\alpha\nr u^2$. \Ex $A\in\mathcal M_q(\R)$ symétrique définie positive : $\sca{Ax}y$, $C=\max\mathrm{Sp}(A)$, $\alpha=\min\mathrm{Sp}(A)$.
\B \Th{Th (Lax-Milgram)} \Hyp $H$ Hilbert, $a$ bilinéaire \textbf{continue et coercive}. \Alors $\forall L\in H'$, $\exists!u\in H$ : $a(u,v)=L(v)$ $\forall v\in H$. \textbf{Si de plus} $a$ symétrique : $u$ est caractérisé par $\frac12a(u,u)-L(u)=\min_{v\in H}\big(\frac12a(v,v)-L(v)\big)$.\\
Preuve : Riesz $\Rightarrow a(u,v)=\sca{Au}v$, $L(v)=\sca fv$ ; $A$ linéaire borné ($\nr{Au}^2=a(u,Au)\le c\nr u\nr{Au}$), injectif ($\sca{Au}u\ge\alpha\nr u^2$), $\mathrm{Im}A$ dense ($\ker A=0$) et fermée ($\alpha\nr{v_n-v_m}^2\le\nr{u_n-u_m}\nr{v_n-v_m}$) : $Au=f$ a une unique solution.
\B \textbf{(D)} $\Omega$ borné régulier (ou $]0,1[$), $c$ continue $\ge0$ : $-\Delta u+cu=f$ dans $\Omega$, $u=0$ sur $\partial\Omega$. Bien posé ? (existence, unicité, continuité par rapport à $f$).
\B \Mt\textbf{Les 5 étapes} : \textbf{(1) formulation faible} : multiplier par $v$ nulle au bord ($C^1$ à support compact), Green : $\int\nabla u\cdot\nabla v+cuv=\int fv$ ; \textbf{(2) formulation variationnelle} : $a(u,v)=\int_\Omega\nabla u\cdot\nabla v+cuv$, $\ell(v)=\int_\Omega fv$, $H=H^1_0$ muni de $(u,v)_{H^1_0}=\int\nabla u\cdot\nabla v$ ; trouver $u\in H$, $a(u,v)=\ell(v)$ $\forall v\in H$ ; \textbf{(3) Lax-Milgram} : coercivité $a(u,u)\ge|u|^2_{H^1}$ (d'où ce produit scalaire : $c$ peut être nulle) ; continuité (Cauchy-Schwarz $+$ Poincaré) $|a(u,v)|\le(1+\nr c_\infty k_\Omega^2)|u|_{H^1}|v|_{H^1}$, $|\ell(v)|\le k_\Omega\nr f_2|v|_{H^1}$ ; donc $\exists!u\in H^1_0$, et $a(u,u)=\ell(u)\Rightarrow|u|_{H^1}\le k_\Omega\nr f_2$ ; \textbf{(4) distributions} : $\mathcal D(\Omega)\subset H^1_0$, $\langle\mathrm{div}\,T,\varphi\rangle=-\langle T,\nabla\varphi\rangle$ $\Rightarrow-\Delta u+cu=f$ dans $\mathcal D'(\Omega)$ ; \textbf{(5) solution classique} : $\Delta u\in L^2$ (distribution régulière), donc $-\Delta u+cu=f$ dans $L^2$. \pg{$\Delta u\in L^2\not\Rightarrow u\in H^2$ en général} ; vrai si bord régulier ou $\Omega$ polygone convexe. En 1D : $u''=cu-f\in L^2$, $u\in H^2\cap H^1_0$, $\nr u_{H^2}\le C\nr f_2$.
\B \Th{Th (résumé)} \Hyp $\Omega$ ouvert borné régulier de $\R^d$, $c\ge0$ \textbf{bornée}, $f\in L^2(\Omega)$. \Alors (D) a une \textbf{unique} solution $u\in H^1_0(\Omega)$ ($-\Delta u+cu=f$ p.p. dans $\Omega$, $u=0$ p.p. sur $\partial\Omega$), et $\nr u_{H^1}\le\mathcal C_\Omega\nr f_{L^2}$ (bien posé dans $H^1$).
\B TD (1D, $f,c\in C^0([0,1])$) : Dirichlet non homogène $u(0)=a$, $u(1)=b$ ($c\ge0$) ; Neumann non homogène $u'(0)=\alpha$, $u'(1)=0$ ($c>0$) ; Dirichlet-Neumann $u(0)=0$, $u'(1)=0$ ($c\ge0$). \Ex $-u''=f$, $u(0)=u(1)=0$ : $f\in L^2\Rightarrow u'\in H^1\Rightarrow u\in H^2\cap H^1_0$ ; $f\ge0\Rightarrow u'$ décroissante et $u\ge0$.

%=====================================================================
\Sec{12. PL continue}
\B $\inf c\cdot x$ s.c. $Ax=b$, $x\ge0$ ; $x\in\R^n$, $A\in\mathcal M_{m\times n}$, $m\le n$, $\mathrm{rg}A=m$. $X_{ad}=\{x\in\R^n_+:Ax=b\}$ polyèdre. $x\in X_{ad}$ : solution \textbf{admissible} ; sommet de $X_{ad}$ : solution \textbf{extrémale}.
\B \Th{Lemme} (1) Le problème admet une solution si l'infimum est minoré \Warn{et $X_{ad}\ne\emptyset$}. (2) Pour $x$ extrémale : permutation des colonnes, $A=[B\ N]$, $B$ inversible, $x=(x_B,x_N)$, $x_N=0$, $x_B=B^{-1}b\ge0$ : solution \textbf{basique}. (3) S'il existe une solution optimale, il en existe une \textbf{basique}. (Idées : ensemble admissible = cône fermé, $x\mapsto c\cdot x$ continue ; un point extrémal n'est pas combinaison convexe non triviale de points de $X_{ad}$.)
\B \Th{Prop} $\bar x$ basique : $\forall x$ admissible, $c\cdot(x-\bar x)=\tilde c_N\cdot x_N$, $\tilde c_N=c_N-{}^t(B^{-1}N)c_B$ (coût réduit) ; $\bar x$ optimale $\Leftrightarrow\tilde c_N\ge0$ \Warn{$\Leftarrow$ toujours ; $\Rightarrow$ si $\bar x_B>0$ (non dégénérée)}.
\B \Th{Rq essentielle} $x$ basique, $\tilde c_N\cdot e_i<0$ : $x(\eps)=\big(B^{-1}(b-\eps N_i),\eps e_i\big)$ vérifie $Ax(\eps)=b$ et $c\cdot(x(\eps)-x)=\eps\,\tilde c_N\cdot e_i<0$ : le coût décroît.

%=====================================================================
\Sec{13. Simplexe et dualité}
\B \Mt\textbf{Simplexe} (1) $x_k$ admissible basique, $A_k=[B_k\ N_k]$ ; $\tilde c_k=c_{k,N}-{}^tN_k\,{}^tB_k^{-1}c_{k,B}$ ; si $\tilde c_k\ge0$ : $x_k$ basique \textbf{optimale} ; sinon $i$ = composante la \textbf{plus négative}. (2) $a_i$ colonne de $N_k$ ; $x_{k,N}(\eps)=\eps e_i$, $x_{k,B}(\eps)=B_k^{-1}(b-\eps a_i)$. (3) Si $\eps$ peut être aussi grand qu'on veut : problème \textbf{non borné} inférieurement. (4) Sinon $\eps_k$ maximal ($x_k(\eps_k)$ admissible, pas au-delà) : une composante $j\ne i$ de $x_{k,B}(\eps_k)$ s'annule ; $x_{k+1}=x_k(\eps_k)$ et on échange les colonnes $i$ et $j$.
\B \textbf{Initialisation} : simplexe pour $\min1\cdot y$ s.c. $\tilde Ax+y=\tilde b$, $(x,y)\ge0$, $\tilde b_i=|b_i|$ ($\tilde A$ correspondant), départ $(y_0,x_0)=(\tilde b,0)$, admissible basique.
\B \textbf{Dualité} : $\Lg(x,p)=c\cdot x+p\cdot(b-Ax)$ ; $G(p)=\min_{x\ge0}\Lg(x,p)=p\cdot b$ si $c-{}^tAp\ge0$, $-\infty$ sinon. Dual : $\sup b\cdot p$ s.c. ${}^tAp\le c$, $p\in\R^m$.
\B \Th{Th} \textbf{Si} le primal ou le dual a une solution finie, \Alors l'autre aussi ; il existe un point selle $(\bar x,\bar p)$ du Lagrangien, admissible pour chacun : $\min c\cdot x=c\cdot\bar x=b\cdot\bar p=\sup b\cdot p$, et \textbf{écarts complémentaires} $\bar x\cdot(c-{}^tA\bar p)=0$.
\B \Ex Centre d'appels : $x_i$ = opérateurs sur la tranche $[\alpha_i,\beta_i]$, salaires $S_i$ : $\min S\cdot x$ s.c. $Ux\ge c$ (clients $c_j$ par intervalle à effectif constant, $u_{ji}=1$ si l'intervalle $j$ est dans la tranche $i$), $x\ge0$ ; $U$ TU.

%=====================================================================
\Sec{14. PL en nombres entiers}
\B $\min_{x\in P\cap\Z^n}c\cdot x$, $P=\{x:Ax\le b\}$ polyèdre : quand vaut-il $\min_Pc\cdot x$ ? \textbf{Enveloppe convexe} = plus petit convexe contenant l'ensemble.
\B \Th{Lemme} $P$ = enveloppe convexe de ses sommets $S$ (points basiques) ; pour $J$ convexe, $\min_PJ=\min_SJ$ \Warn{faux pour $J$ convexe ; vrai pour $J$ concave, en particulier linéaire, et $P$ borné}.
\B \Th{Lemme} $A\in\mathcal M_n(\Z)$ inversible : ($\forall b\in\Z^n$, $A^{-1}b\in\Z^n$) $\Leftrightarrow\det A=\pm1$ (\textbf{unimodulaire}) ; alors, si $b\in\Z^n$, $P$ = enveloppe convexe de $P\cap\Z^n$.
\B \textbf{Totalement unimodulaire} ($m\ne n$) : toute sous-matrice carrée est unimodulaire ou non inversible. \Th{Lemme} Coefficients dans $\{-1,0,1\}$, au plus un $1$ et un $-1$ par colonne $\Rightarrow A$ TU.
\B \Th{Corollaire} \Hyp $A$ TU, $b\in\Z^n$ ($c\in\R^n$). \Alors les solutions basiques de $\min_Pc\cdot x$ sont \textbf{entières} ; si un min existe, il est atteint en une solution basique entière.

%=====================================================================
\Sec{15. Flots}
\B Graphe orienté $G=(\mathcal N,\mathcal A)$, $\mathcal A\subset\mathcal N\times\mathcal N$, $m$ nœuds, $n$ arcs ; arc $(i,j)$ de $i$ vers $j$, capacité $u_{ij}\ge0$, coût $c_{ij}$, flot extérieur $b_i$ (entrant ou sortant).
\B \textbf{Flot admissible} : $x\in\R^n_+$, loi des nœuds (Kirchhoff) $b_i+\sum_{(j,i)\in\mathcal A}x_{ji}=\sum_{(i,j)\in\mathcal A}x_{ij}$ $\forall i$, capacités $x_{ij}\le u_{ij}$.
\B \Th{Propriété} $A\in\mathcal M_{m\times n}$, $A_{i,l}=1$ si $l=(j,i)$, $-1$ si $l=(i,j)$, $0$ sinon : la loi des nœuds s'écrit $Ax=b$ \Warn{avec ces signes c'est $Ax=-b$} ; $A$ est \textbf{TU}.
\B \textbf{Flot de coût minimum} : $\min c\cdot x$ s.c. $0\le x\le u$, $Ax=b$ (affectations, emplois du temps, approvisionnement) ; existence si $c_{ij}\ge0$ ou $u_{ij}<\infty$ ; \textbf{si $u$ et $b$ entiers}, le simplexe fournit une solution optimale \textbf{entière}. \Ex 2 usines (100, 80) vers 2 concessionnaires (80, 75) : ajouter un concessionnaire fictif qui reçoit le stock restant à coût nul.
\B \textbf{Flot maximal} (source $s$ sans antécédent, puits $p$ sans successeur) : flot admissible maximisant $b_s$, avec $b_p=-b_s$ (conservation globale), $b_i=0$ sinon ; on ajoute l'arc $(p,s)$ et on augmente $x_{ps}$ depuis 0 tant que c'est admissible ; les $c_{ij}$ n'interviennent pas.
\B \textbf{Chemin de coût minimum} : flot de coût min avec $x_{ij}\in\{0,1\}$, $b_s=-b_p=1$, $b_i=0$ sinon ; solutions optimales connexes (si $x_{ji}=1$ alors $x_{ij'}=1$ pour un $j'\ne j$) ; le simplexe résout la relaxation $0\le x_{ij}\le1$, mais \textbf{Dijkstra} est meilleur (complexité).

%=====================================================================
\Sec{16. Problèmes entiers célèbres}
\B \textbf{Sac à dos} : $N$ objets, poids $p_j$, utilités $u_j$ : $\max u\cdot x$, $x_j\in\{0,1\}$, $p\cdot x\le L$. \textbf{NP-complet} (complexité exponentielle) ; heuristique : algorithme \textbf{glouton} ; résolution exacte par Ford-Bellman.
\B \textbf{Programmation dynamique} : $D=c_{x_0x_1}+\dots+c_{x_{N-1}x_N}+b_{x_N}$ (coût ou revenu d'un processus d'états $x_i$, graphe des transitions) ; horizon fini $N$, chemins issus de $x_0$. \textbf{Principe de Bellman} : tout sous-chemin d'un chemin optimal est optimal.
\B \Th{Prop} Fonction valeur $v^N:x_0\mapsto v^N_{x_0}$ (solution à horizon $N$). Pour le min : $v^0_x=b_x$, $v^n_x=\inf_{y\in\mathcal N}\big(c_{xy}+v^{n-1}_y\big)$, $1\le n\le N$.
\B Sac à dos : états = poids $q\in[0,L]$, $v^n_q$ = utilité d'un sac partiel rempli de $n$ objets : $v^n_q=\max\big(v^{n-1}_q,\max_i(u_i+v^{n-1}_{q-p_i})\big)$ \Warn{objet $n$ pris au plus une fois : $\max(v^{n-1}_q,u_n+v^{n-1}_{q-p_n})$}.
\B \textbf{Voyageur de commerce} : $\min_{\sigma\in\mathcal E_N}\sum d(x_{\sigma(n)},x_{\sigma(n+1)})$ ; binaire $x\in\{0,1\}^{\mathcal A}$ : (1) $\sum_jx_{ij}=2$ $\forall i$ ; (2) connexité $\sum_{(i,j)\in S\cap\mathcal A}x_{ij}\le\mathrm{card}(S)-1$ $\forall S$ ; (3) $\min\sum d_{ij}x_{ij}$ s.c. (1), (2).
\B \textbf{Séparation-évaluation} (branch and bound) : arbre des tours ; borne $m$ initialisée par une heuristique ; borne inférieure prévisionnelle par relaxation des contraintes et dualité (relaxation lagrangienne) ; on élague les branches dont la borne prévisionnelle excède $m$.
\B \textbf{Liens logiques} (binaires $x_a,x_b,x_c$) : pas $a$ et $b$ ensemble : $x_a+x_b\le1$ ; au moins un : $x_a+x_b\ge1$ ; $a\Rightarrow b$ : $x_a\le x_b$ (contient la contraposée) ; $a\Rightarrow(b$ et/ou $c)$ : $x_a\le x_b+x_c$.

\end{multicols}

\end{document}
